\begin{itemize}
\item \textbf{Simulated, i.i.d.\ NK only.} Contributions are i.i.d.\ uniform tables with random neighbours; real proteins have correlated effects, global epistasis, sparsity, noise and structured neighbourhoods. The results say how these models behave on this family and nothing about proteins.
\item \textbf{Scale.} Eight tasks at $N=1000$, five seeds, two $(L,A)$ pairs, $K\in\{0,1,2,4\}$; $K=3$, intermediate alphabets and noisy labels were not run. $L$ and $A$ change together between the two settings, so no difference between them can be attributed to alphabet size or to length alone. The $N$ sweep covers two tasks and none of the regimes where every model fails. Per-cell differences with five seeds are often within seed variation.
\item \textbf{Reimplemented, lightly and unequally tuned baselines.} The CNN and MLP have one architecture each and 2 configurations (weight decay) against 24 for pairwise ridge and 6 for additive ridge, and they are fit on 85\% of the training set without the refit on all data that the ridge models get. Both asymmetries favour ridge; larger networks, a refit, ensembles, other regularisers or pretraining could change the net-versus-ridge comparison (H3).
\item \textbf{Test distribution.} Random train and test sequences; no homology, no extrapolation split, no noise.
\item \textbf{Design endpoint.} One greedy step of 10 proposals from one parent; \emph{design\_hit} is quantised in tenths per run, so with five seeds only large differences are detectable: outside $K=0$ and pairwise ridge at L15\_A4\_K1, the design\_hit differences in this paper are at most nominally significant in uncorrected tests, and the H4 correlation rests on that one cell. The parent is the best training sequence (selected on noiseless fitness), and no closed-loop or batch-diversity design was tried.
\end{itemize}
